\section{Introduzione}
\label{sec:introduzione}

L'autenticazione digitale costituisce il primo livello di difesa per proteggere l'accesso agli account personali e alle risorse informatiche. Con la crescita dei servizi online, i meccanismi di verifica dell'identità sono passati da credenziali gestite singolarmente da ogni sito a sistemi centralizzati basati su un unico account.

\subsection{Limiti delle password e nascita del login federato}
Per lungo tempo l'autenticazione sul web si è basata quasi esclusivamente sulla coppia \emph{username} e \emph{password} gestita localmente da ciascun servizio. Con la moltiplicazione delle piattaforme utilizzate quotidianamente, questo approccio ha mostrato limiti evidenti di usabilità e sicurezza. Ricordare decine di credenziali complesse induce spesso gli utenti a scegliere password deboli o a riutilizzarle su più siti~\cite{nist800-63b}: la compromissione di un singolo servizio tramite \emph{data breach} espone così gli account dell'utente su altre piattaforme a campagne automatizzate di \emph{credential stuffing}. Inoltre, per i gestori delle applicazioni, custodire internamente i database delle credenziali comporta costi e responsabilità rilevanti.

Per superare questa frammentazione si è affermato il modello del \textbf{login federato}. Nato in ambito aziendale e accademico (ad esempio con lo standard SAML o infrastrutture come \texttt{eduroam}), si è poi diffuso nel web consumer attraverso i \emph{social login} come «Accedi con Google», «Accedi con Facebook» e «Sign in with Apple»~\cite{nist800-63c}. Il login federato consente di sfruttare un'identità digitale già esistente presso un fornitore principale per autenticarsi su applicazioni di terze parti senza creare nuove password, migliorando l'usabilità (soprattutto su dispositivi mobili) e delegando la sicurezza a provider dotati di infrastrutture avanzate e autenticazione a più fattori (MFA). Tuttavia, concentrare l'autenticazione su pochi fornitori sposta la superficie d'attacco sulla corretta verifica delle attestazioni scambiate tra i sistemi coinvolti, come dimostra la vulnerabilità scoperta nel 2020 in \emph{Sign in with Apple} analizzata in questa relazione.